\section{程序总体思路}

本程序对应的源文件是
\[
\texttt{1-D\_Riemann\_NM\_MacC\_02.c}.
\]
它求解一维无粘可压 Euler 方程：
\[
\frac{\partial \bm Q}{\partial t}
+\frac{\partial \bm F(\bm Q)}{\partial x}=0,
\qquad
\bm Q=(\rho,\rho u,\rho E)^{\mathrm T}.
\]

这份代码主要用于 Riemann 初值问题，例如 Sod shock tube。初始时刻左右两边给定不同的
\((\rho,u,p)\)，中间在 \(x_0\) 处有一个间断。程序把左右原始量转换成守恒量，
然后用 MacCormack 格式向前推进。

\subsection{一个时间步的大致流程}

02 版使用的顺序是：
\[
\bm Q^n
\rightarrow
\bar{\bm Q}^{\,n}
\rightarrow
\tilde{\bm Q}^{\,n+1}
\rightarrow
\bm Q^{n+1}.
\]

这里：
\begin{itemize}
    \item \(\bm Q^n\)：当前时刻的守恒量；
    \item \(\bar{\bm Q}^{\,n}\)：加人工粘性滤波后的状态；
    \item \(\tilde{\bm Q}^{\,n+1}\)：预测步得到的状态；
    \item \(\bm Q^{n+1}\)：校正步后的新状态。
\end{itemize}

程序中的主循环基本就是反复调用：
\begin{lstlisting}
advance_one_maccormack_step(&solver);
\end{lstlisting}
每调用一次，就完成一个时间步。

\subsection{主程序做了什么}

\texttt{main()} 的执行顺序如下：

\begin{enumerate}[label=\arabic*.]
    \item 建立一个带默认参数的 \texttt{Solver}。
    \item 选择 Riemann 算例，例如 Sod、Lax 或纯接触间断。
    \item 输入计算域、网格数、CFL、终止时间和人工粘性参数。
    \item 给所有数组分配内存。
    \item 初始化守恒量 \(\rho,\rho u,\rho E\)。
    \item 进入时间推进循环。
    \item 输出 Tecplot 数据，并调用 Project 0 生成精确解。
\end{enumerate}

因此，这份代码最主要的整理工作是把 CFD 变量按物理意义和时间层次放好。
这样后面检查程序时，不需要反复把 \texttt{q1/q2/q3} 和具体物理量对应起来。
